\section{Introduction}
\label{sec:introduction}
Python's dynamic type system offers developers considerable flexibility
and enables rapid development. However, it also allows type errors to
remain undetected until runtime. These errors can interrupt program
execution and cause application failures, compromising software
reliability. Consequently, substantial research effort has been devoted
to static type-error detection for Python, with the aim of identifying
such errors before execution.

Existing approaches have made progress in static type-error detection
for Python. Type checkers such as Mypy and Pyrefly use type annotations
to check assignments, function calls, and other operations against declared
types, and have been adopted in major production codebases and open-source
projects, including Dropbox, Instagram, PyTorch, and
JAX~\cite{dropboxMypy,pyreflyAdoption}. Automatic type inference further
improves detection. For example, Pyinder detected 50\% (34/68) of real-world type
errors  in its benchmark~\cite{oh2024pyinder}.
However, these approaches rely on static type models that do not fully
capture how types vary with execution conditions or change through function
calls. As a result, they can produce false alarms or miss type errors.
In Pyinder's evaluation, such limitations accounted for 31\% of the
confirmed false alarms in its inspected sample~\cite{oh2024pyinder}.
This gap between static type models and Python's dynamic typing behavior
limits the effectiveness of existing type-error detection methods.

To address this gap, we aim to improve Python type-error detection by
tracking type changes across function calls. A function's type requirements
concern not only its arguments but also the global variables, instance
attributes, and other values it accesses. We refer to their collective type
information at a given program point as the \emph{type-state}. Our approach
captures how a function transforms this state under different execution
conditions and uses these transformations to check subsequent operations.
To achieve this goal, we need to address three key challenges.

\paragraph{Challenge 1: Modeling type-state changes across function calls.}
A function can change the caller's type-state not only through its return
value but also by updating global variables or attributes of mutable
objects. These updates depend on the function's execution path, so calls
to the same function can establish different type-states. Conventional
type annotations describe parameter, return, and attribute types, but do
not specify how a call transforms the incoming type-state under different
execution conditions. We therefore need a new model to express these
conditional type-state changes and their effects on subsequent operations.

\paragraph{Challenge 2: Extracting conditional type-state changes from code.}
Program execution involves both explicit and implicit paths. Explicit
paths arise from branch statements, while implicit paths arise from Python
operations whose behavior depends on the types involved. For example,
addition can perform arithmetic, concatenate strings, or invoke a
user-defined method, with different type requirements and effects in each
case. Analyzing these paths across function calls further requires
considering combinations of caller and callee paths, leading to path
explosion. Accounting for implicit paths and managing path explosion
therefore make it challenging to extract conditional type-state changes
from code.

\paragraph{Challenge 3: Detecting type errors using conditional type-state changes.}
Due to limitations in analysis precision, the extracted type-state
transition relations may contain inaccuracies, which can cause valid
operations to be flagged as type errors. A detected type conflict therefore
does not necessarily indicate
an error in the program. Distinguishing a real error from an analysis
inaccuracy requires examining how the conflicting state arises under the
relevant calling conditions. Moreover, an inaccurate type-state change
can propagate through multiple calls, affecting error detection in other
functions. The challenge is therefore to distinguish type errors
from false alarms and correct analysis inaccuracies.

To address these challenges, we propose PyTT(Python Type Tracker), a Python type-error detector
that models and tracks conditional type-state transitions across function
calls. To address \textbf{Challenge 1}, we introduce a
\emph{path-sensitive type contract} (PSTC) to represent these transitions,
capturing a function's type requirements and updates along different
execution paths. To address \textbf{Challenge 2}, PyTT uses static analysis
to identify explicit paths and a large language model (LLM) to infer
implicit paths. It further uses static analysis to
simplify code whose type behavior is independent of the calling context,
reducing repeated analysis across paths to mitigate path explosion.
To address \textbf{Challenge 3}, PyTT uses PSTCs to trace detected conflicts
across function boundaries. It analyzes the resulting traces to distinguish
type errors from false alarms. For false alarms, PyTT corrects the
inaccuracies in the corresponding contracts and updates the affected
analysis results.

We evaluate PyTT on xxx real-world type-error fixes from xxx Python
repositories, measuring detection effectiveness, diagnostic precision,
and analysis cost. PyTT detects xxx runtime type errors and xxx
annotation-contract errors, with precision of xxx\% and xxx\%, respectively.
Each detection is checked against both the buggy and fixed revisions.
Ablation experiments examine the contributions of conditional type-state
modeling, explicit and implicit path analysis, and conflict analysis and
correction. We further evaluate PyTT on xxx independent repositories,
where it identifies xxx previously unknown type errors.

The contributions are as follows:
\begin{itemize}
  \item We introduce the path-sensitive type contract (PSTC), a model
  that captures conditional type-state transitions across function calls
  for Python type-error detection.
  \item We propose PyTT, a type-error detection method that combines
  static analysis and LLM-based analysis to extract these transitions.
  PyTT analyzes conflict traces across functions to distinguish type
  errors from false alarms and correct analysis inaccuracies.
  \item We evaluate PyTT on real-world type-error fixes and independent
  Python repositories, assessing detection effectiveness, diagnostic
  precision, and analysis cost, and examining the contributions of its
  key components through ablation experiments.
\end{itemize}
